\section{Methodology}\label{methodology}

\subsection{\texorpdfstring{\textbf{Research Design}}{Research Design}}\label{research-design}

This study used a quantitative, non-experimental design to model residential property values in Metro Cebu. The analysis was predictive in that it estimated prices from observed property characteristics, and prescriptive in that the results were later organized for spatial decision support through GIS. The training target across all three models was total property price in Philippine Peso, log-transformed as \emph{log\_price} for model fitting. Price per square meter (\emph{price\_per\_sqm}) served as the substantive unit-price quantity of interest and as a derived diagnostic field; evaluation results were back-transformed to total price in PHP for reporting. The independent variables included structural, geospatial, and administrative features. Macroeconomic indicators were reviewed during study design but were not retained in the final fitted specification.

Three supervised learning approaches were compared in order to balance interpretability and predictive performance, while also testing whether GIS-derived and administrative variables improved the model beyond structural features alone:

\begin{enumerate}
\item \textbf{Ordinary Least Squares (OLS) / Hedonic Regression}: An interpretable baseline grounded in Hedonic Pricing Theory (Rosen, 1974). Coefficients carried direct economic meaning (e.g., ``each additional bedroom adds PHP X to value'').
\item \textbf{Random Forest Regressor}: A tree-based ensemble that captured non-linear relationships and feature interactions without requiring explicit specification (Breiman, 2001).
\item \textbf{XGBoost Regressor}: A gradient boosting algorithm optimized for predictive accuracy on structured tabular data (Chen \& Guestrin, 2016). Empirically, XGBoost often performed well on datasets of similar scale (Nyanda et al., 2024).
\end{enumerate}

\begin{center}\rule{0.5\linewidth}{0.5pt}\end{center}

\subsection{\texorpdfstring{\textbf{Data Sources: Multi-Source Property Evidence}}{Data Sources: Multi-Source Property Evidence}}\label{data-sources-the-hybrid-strategy}

Because direct deed-of-sale transaction records were not publicly accessible, the broader property evidence base was assembled from multiple sources that reflected different segments of the residential market. Foreclosure and acquired-asset records were retained as conservative benchmark layers, while open-market listings provided the labels used for the deployed prediction task. The purpose of this multi-source design was not to claim direct observation of a single true market price, but to preserve the wider valuation context while keeping the final fitted model aligned with open-market residential pricing.

\begin{table}[htbp]
\caption{Summary of Data Sources Used in the Thesis Workflow}
\label{tab:chapter3-data-sources}
\begin{center}
\begin{tabular}{p{0.26\textwidth} p{0.20\textwidth} p{0.22\textwidth} p{0.20\textwidth}}
\toprule
Source & Role & Volume & Nature \\
\midrule
\textbf{Open-Market Listings} & Deployed training label source & Lamudi residential listings & Asking / market-facing \\
\textbf{Bank Foreclosure Inventory} & Conservative benchmark layer & BDO, BPI, Metrobank, Landbank, Bank of Commerce, China Bank Savings & Conservative / distressed \\
\textbf{Floor-Price References} & Conservative benchmark layer & Pag-IBIG and BDO records & Administrative / distressed \\
\textbf{BIR Zonal Values} & Administrative benchmark & Per barangay & Static / regulatory \\
\textbf{Google Maps APIs} & Geocoding and POI retrieval & Per property & Geospatial / dynamic \\
\textbf{OpenStreetMap Road Network} & Routing and network topology & Per property / destination pair & Geospatial / open data \\
\bottomrule
\end{tabular}
\end{center}
\end{table}

\subsubsection{\texorpdfstring{\textbf{Bank Foreclosure and Floor-Price Records}}{Bank Foreclosure and Floor-Price Records}}\label{primary-data-verified-floor-price-foreclosed-acquired-assets}

The conservative valuation layers aggregated foreclosed and acquired property listings from multiple institutional sources to prevent single-source bias:

\begin{itemize}
\item \textbf{BDO Unibank}: The largest single foreclosure source used in the study and the main template for raw-to-clean pipeline design.
\item \textbf{Pag-IBIG Fund (HDMF)}: Acquired assets and floor-price references representing the more affordable residential segment.
\item \textbf{Other bank foreclosure sources}: BPI, Metrobank, Landbank, Bank of Commerce, and China Bank Savings records were added to widen conservative market coverage.
\end{itemize}

These distressed assets were treated as a lower-bound segment of the observed market rather than as direct equivalents of open-market transaction prices. Using multiple institutional sources helped reduce the bias that might result if the dataset were drawn from only one lender or seller.

A sample of the raw BDO and bank ROPA input structure is documented in Appendix A.

\subsubsection{\texorpdfstring{\textbf{Open-Market Listings}}{Open-Market Listings}}\label{secondary-data-market-ceiling-price-online-listings}

To represent current asking prices in the residential market, the study collected listings from public online platforms, primarily Lamudi. As Sousa et al. (2024) note, large listing datasets can capture pricing patterns that are often missing from sparse official records. In the final thesis workflow, these open-market listings became the training scope for the deployed price-surface model.

A sample of the cleaned Lamudi listing structure is documented in Appendix A.

\subsubsection{\texorpdfstring{\textbf{Administrative and Macroeconomic Data}}{Administrative and Macroeconomic Data}}\label{administrative-and-macroeconomic-data}

\begin{itemize}
\item \textbf{BIR Zonal Values}: Official zonal values per barangay, used as administrative benchmark features and as the basis for the derived valuation-gap diagnostic.
\item \textbf{RPPI Review}: The BSP residential price index was reviewed during study design as a possible macro-control, but it was not retained in the final fitted specification. The deployed model is cross-sectional and listing-level, so quarterly RPPI movements were not joined as an active model feature in the final ABT.
\end{itemize}

\subsubsection{\texorpdfstring{\textbf{Geospatial Data Sources}}{Geospatial Data Sources}}\label{geospatial-data-sources}

\begin{itemize}
\item \textbf{Google Maps Geocoding API}: Converted property addresses into latitude/longitude coordinates and supported reverse geocoding of barangay names for the BIR join.
\item \textbf{Google Maps Places API}: Supplied the raw point-of-interest inventories that were later curated into the nine MCRAI amenity categories.
\item \textbf{OpenStreetMap Road Network}: Provided the street graph used for network-distance computation through the \emph{osmnx} workflow \parencite{boeing2017osmnx,boeing2019urban}. In the final implementation, OSM was used for routing and network topology rather than as the primary POI source.
\end{itemize}

\subsubsection{\texorpdfstring{\textbf{Target Variable}}{Target Variable}}\label{target-variable}

The canonical ABT retained three market tiers---\emph{open\_market}, \emph{bank\_ropa}, and \emph{floor\_price}---to preserve the broader valuation evidence base. However, the predictive task of the thesis is the estimation of \emph{open-market residential price}. Accordingly, model fitting filtered the ABT to \emph{market\_segment = open\_market} before estimation, while foreclosure and administrative floor-price records were kept as reference layers in the wider data system rather than as training labels for the deployed model. The primary target variable for model comparison was \emph{price\_per\_sqm}; total price and its log transform were retained as derived fields for diagnostics and the OLS baseline.

The complete feature matrix schema is documented in Appendix A.
\subsubsection{\texorpdfstring{\textbf{Validation Layer: Human-in-the-Loop}}{Validation Layer: Human-in-the-Loop}}\label{validation-layer-human-in-the-loop}

Licensed real estate brokers from the CPRE network were included as a validation layer, not to replace statistical evaluation, but to check whether the model outputs remained plausible within local market practice. Their role was limited to two tasks:

\begin{enumerate}
\item \textbf{Sanity Check}: Reviewing SHAP-derived value driver rankings against practitioner knowledge.
\item \textbf{Outlier Review}: Examining cases with high prediction error to distinguish likely data issues from genuine market anomalies.
\end{enumerate}

\begin{center}\rule{0.5\linewidth}{0.5pt}\end{center}

\subsection{\texorpdfstring{\textbf{Data Pipeline}}{Data Pipeline}}\label{data-pipeline}

The data preparation process proceeded in seven stages, implemented in Python.

\begin{enumerate}
\item \textbf{Ingestion}: BDO acquired asset data was ingested from Excel via Pandas. Lamudi listings were collected through a custom web scraper. Pag-IBIG foreclosed property records were ingested from structured PDFs and Excel exports.
\item \textbf{Filtering}: The dataset was restricted to residential properties within Metro Cebu---Cebu City, Mandaue City, Lapu-Lapu City, Talisay City, Minglanilla, and Consolacion. The canonical ABT retained all three market tiers (\emph{open\_market}, \emph{bank\_ropa}, and \emph{floor\_price}), but the later modeling scripts filtered this ABT to \emph{open\_market} observations only before fitting.
\item \textbf{Regex Parsing and Cleaning}:
\begin{itemize}
\item \emph{BDO Data}: The \emph{Property Description} field bundled features (e.g., ``3BR 2TB''). Regex patterns extracted \emph{Bedrooms} and \emph{Bathrooms}.
\item \emph{Lamudi Data}: Scraped price fields contained varying formats (e.g., ``PHP 5,000,000'' or ``Contact agent for price''). Currency symbols and commas were stripped; rows without explicit numerical prices were dropped.
\item \emph{Pag-IBIG Data}: Minimum bid prices were extracted and tagged as floor-price observations.
\end{itemize}
\item \textbf{Geocoding (Address to Coordinates)}: Property addresses were batch-geocoded through the Google Maps Geocoding API to obtain latitude and longitude. The service was selected for project-level address resolution and reverse geocoding in Metro Cebu, where listings often reference barangays, landmarks, or compound descriptors rather than standardized street-number formats. Results were cached locally to avoid redundant API calls.
\item \textbf{BIR Zonal Value Extraction and Barangay Join}: Official BIR zonal value schedules were sourced from four Revenue District Office files covering the six-LGU study area. A custom parser extracted street-level values by residential classification code and aggregated them to the barangay level as the median per classification. Each property was then assigned a barangay through reverse geocoding via the Google Maps Geocoding API, which served as the join key to the BIR summary table. This process yielded full bir\_zonal\_rr\_median coverage across all 2{,}047 ABT rows.
\item \textbf{GIS Augmentation}: From geocoded coordinates, the geospatial feature set described in \S3.4.1 was computed: road-network distances to eight polycentric and regional anchor nodes; nine category-specific MCRAI component scores derived from curated amenity inventories; the \emph{mcrai\_composite}; and a spatial lag variable capturing the mean price of neighboring properties within 1 km. Network shortest paths were computed on the OSM road graph, with Haversine fallback only when a point could not be snapped cleanly to the network.
\item \textbf{Final ABT Assembly}: All engineered features were merged into a single Analytics Base Table (ABT) and saved as a flat CSV for modeling. The canonical ABT preserved all market tiers, while the modeling-ready subset was produced downstream by filtering to \emph{open\_market}, dropping rows with null \emph{price\_per\_sqm} or \emph{spatial\_lag\_price}, and excluding three integrity-flagged price outliers.
\end{enumerate}

\textbf{Tools}: Python (Pandas, NumPy, Scikit-learn, XGBoost, Statsmodels, Requests, OSMnx), Google Maps Geocoding and Places APIs, and QGIS for spatial visualization.

\begin{center}\rule{0.5\linewidth}{0.5pt}\end{center}

\subsection{\texorpdfstring{\textbf{Feature Engineering}}{Feature Engineering}}\label{feature-engineering}

\begin{table}[htbp]
\caption{Feature Categories Used in the Thesis Workflow}
\label{tab:chapter3-feature-categories}
\begin{center}
\begin{tabular}{p{0.18\textwidth} p{0.58\textwidth} p{0.16\textwidth}}
\toprule
Category & Features & Source \\
\midrule
\textbf{Structural} & Lot Area, Floor Area, Bedrooms, Bathrooms, Property Type, vacancy / imputation flags & BDO / Lamudi / bank records \\
\textbf{Locational} & City, Barangay, Latitude/Longitude, \emph{is\_mactan\_island} & Google Maps API \\
\textbf{Geospatial} & Network distance to Cebu Business Park, Mandaue CBD, Mactan CBD, SRP, Talisay Tabunok, Consolacion, Naga City, and Airport & Geocoding + OSM road network \\
\textbf{MCRAI} & Nine component scores (education, finance, grocery, health, security, transport, tourism, recreation, retail density) plus \emph{mcrai\_composite} & Places API POI inventory + network routing \\
\textbf{Spatial Lag} & Mean price of neighboring properties within 1 km & Computed from data \\
\textbf{Administrative} & BIR zonal-value features per barangay & BIR schedules \\
\textbf{Metadata} & Source and market-segment fields retained in ABT but excluded from the final open-market feature matrix & Engineered \\
\bottomrule
\end{tabular}
\end{center}
\end{table}

\textbf{Engineered variables}: \emph{price\_per\_sqm}, \emph{valuation\_gap} (a diagnostic field, not a model feature), and derived log-price terms for the baseline regression workflow.

\subsubsection{\texorpdfstring{\textbf{Geospatial Feature Engineering}}{Geospatial Feature Engineering}}\label{geospatial-feature-engineering}

Geospatial feature construction was the part of the method that most directly connected the valuation model to the urban structure of Metro Cebu. Rather than treating location as a background descriptor, this stage translated network accessibility, amenity access, and neighborhood context into variables that could enter the model explicitly.

\begin{enumerate}
\item \textbf{Geocoding and Administrative Join}: Property addresses were geocoded through the Google Maps Geocoding API to obtain latitude and longitude coordinates, then reverse-geocoded to recover barangay names for the BIR zonal-value join described in \S3.3. Geocoding results were cached locally to preserve API quota and ensure repeatable downstream feature construction.

\item \textbf{Network Distance Features}: For each property, shortest-path road distances were computed to eight polycentric and regional anchor nodes: Cebu Business Park, Mandaue CBD, Mactan CBD, South Road Properties, Talisay Tabunok, Consolacion, Naga City, and Mactan-Cebu International Airport. This node set followed the thesis' polycentric framing of Metro Cebu and retained both internal subcenters and major external anchors that shape residential accessibility gradients \parencite{giuliano1991subcenters}. Distances were computed on an OpenStreetMap street graph using the OSMnx workflow \parencite{boeing2017osmnx,boeing2019urban}. When a property or destination could not be snapped cleanly to the network, a Haversine fallback was used so that valid observations were not dropped from the ABT.

\item \textbf{Metro Cebu Residential Accessibility Index (MCRAI)}: Amenity access was operationalized through the Metro Cebu Residential Accessibility Index (MCRAI), a custom gravity-style accessibility framework adapted to local valuation practice. Category-specific point inventories were assembled from Google Maps Places queries and cleaned into nine categories: education, finance, grocery, health, security, transport, tourism, recreation, and retail density. Following Hansen-style accessibility logic, nearer opportunities contribute more strongly than farther opportunities \parencite{hansen1959accessibility}. For property \(i\) and category \(c\), the category score was computed as:

\[
MCRAI_{ic} = \sum_{j \in c} \frac{1}{\max(d_{ij}, 0.5)^{\beta}}, \qquad \beta = 2,
\]

where \(d_{ij}\) is the road-network distance in kilometers from property \(i\) to amenity \(j\), floored at 0.5 km to prevent singular values for immediately adjacent points. Each category was evaluated only within its designated search radius \(r_c\), reflecting the fact that some amenities operate at neighborhood scale while others serve wider urban catchments.

\begin{table}[htbp]
\caption{MCRAI Categories and Search Radii}
\label{tab:chapter3-mcrai-radii}
\begin{center}
\begin{tabular}{p{0.20\textwidth} p{0.34\textwidth} p{0.12\textwidth}}
\toprule
Category & Representative POIs & Radius (km) \\
\midrule
Education & Schools, universities & 0.8 \\
Finance & Banks, ATMs & 1.5 \\
Grocery & Supermarkets, wet markets & 2.0 \\
Health & Hospitals, clinics & 2.0 \\
Security & Police substations, fire stations, barangay halls & 2.0 \\
Transport & Road-corridor / highway way centers & 3.0 \\
Tourism & Hotels, resorts, transient lodging & 3.0 \\
Recreation & Gyms, sports complexes, parks, beaches & 1.5 \\
Retail Density & Restaurants, cafes, bakeries, convenience stores & 1.0 \\
\bottomrule
\end{tabular}
\end{center}
\end{table}

The nine \emph{mcrai\_*} component scores entered the model individually. A separate \emph{mcrai\_composite} was retained as a summary index, but its final specification retained only the positive-coefficient OLS categories---education, grocery, recreation, and transport---in the composite, with weights 0.401, 0.310, 0.199, and 0.102 respectively. Security, tourism, retail density, finance, and health remained available as stand-alone feature columns, but they were not used in the composite weight vector. For the negative-coefficient categories in particular, the preferred interpretation is spatial sorting or threshold externality rather than direct residential amenity value generation \parencite{tiebout1956,bayer2012tiebout,dronyk2017goldilocks,yang2016externality,song2004mixed,chenjim2010}.

\item \textbf{Spatial Lag Variable}: To capture neighborhood price effects, \emph{spatial\_lag\_price} was computed as the mean price of other properties within a 1 km neighborhood of each target property. Properties with no neighbors within this radius received a null lag and were excluded from model fitting. This variable operationalized Tobler's First Law of Geography in a form that could be used directly by the non-spatial regression and tree models \parencite{tobler1970movie}.
\end{enumerate}
\begin{center}\rule{0.5\linewidth}{0.5pt}\end{center}

\subsection{\texorpdfstring{\textbf{Pre-processing}}{Pre-processing}}\label{pre-processing}

\begin{table}[htbp]
\caption{Pre-processing Steps Applied to the Modeling Subset}
\label{tab:chapter3-preprocessing}
\begin{center}
\begin{tabular}{p{0.24\textwidth} p{0.28\textwidth} p{0.34\textwidth}}
\toprule
Step & Method & Rationale \\
\midrule
\textbf{Integrity Filtering} & Remove rows with null \emph{price\_per\_sqm}, null \emph{spatial\_lag\_price}, or explicit integrity flags & Keeps the modeling subset consistent with the final pipeline \\
\textbf{Log Transformation} & Apply log transform to the price target used by the OLS baseline & Reduces right skew while preserving the open-market target definition \\
\textbf{Missing Values} & Median imputation for remaining numeric fields; explicit zeros for structurally absent or all-null fields & Preserves usable rows while handling sparse engineered fields \\
\textbf{Encoding} & Dummy encoding for city and property type & Required for the fitted design matrix while keeping barangay for joins and reporting \\
\bottomrule
\end{tabular}
\end{center}
\end{table}

\begin{center}\rule{0.5\linewidth}{0.5pt}\end{center}

\subsection{\texorpdfstring{\textbf{Modeling Strategy}}{Modeling Strategy}}\label{modeling-strategy}

\subsubsection{\texorpdfstring{\textbf{The Three Models}}{The Three Models}}\label{the-three-models}

The rationale for selecting these three models, including alternatives that were considered but excluded, is described in Section 1.6.

\subsubsection{\texorpdfstring{\textbf{Hedonic Equation}}{Hedonic Equation}}\label{hedonic-equation}

The hedonic regression took the following log-linear form for the baseline comparator:

\[
\begin{aligned}
\ln(\text{price\_php}_i) ={}& \alpha + \beta_1 \ln(\text{floor area}_i) + \beta_2(\text{bedrooms}_i) + \beta_3(\text{bathrooms}_i) \\
& + \beta_4(\text{distance-to-node features}_i) + \beta_5(\text{MCRAI components}_i) \\
& + \beta_6(\text{BIR zonal features}_i) + \beta_7(\text{spatial lag}_i) + \epsilon_i
\end{aligned}
\]

The dependent variable is the log of total property price in Philippine Peso (\emph{log\_price}), consistent with the training target used for OLS, Random Forest, and XGBoost. Predicted values are back-transformed to total price PHP before evaluation metrics are computed. In this form, the hedonic model retains its interpretive structure while allowing administrative and spatial variables to enter directly into the price equation. This makes it possible to compare a more conventional valuation framework with the tree-based models without stripping out the locational features that are central to the study.

\subsubsection{\texorpdfstring{\textbf{Hyperparameter Tuning}}{Hyperparameter Tuning}}\label{hyperparameter-tuning}

For the ML models (Random Forest and XGBoost), exploratory tuning was run in a separate script using \textbf{RandomizedSearchCV} with \textbf{5-fold cross-validation}. Candidate parameter combinations were sampled for both algorithms and compared against the untuned baseline specifications. The baseline Random Forest and XGBoost models were retained for the main pipeline because the tuned variants did not deliver sufficient improvement to replace the simpler benchmark configuration.

\begin{center}\rule{0.5\linewidth}{0.5pt}\end{center}

\subsection{\texorpdfstring{\textbf{Evaluation and Explainability}}{Evaluation and Explainability}}\label{evaluation-and-explainability}

\subsubsection{\texorpdfstring{\textbf{Performance Metrics}}{Performance Metrics}}\label{performance-metrics}

\begin{table}[htbp]
\caption{Performance Metrics Used in Model Evaluation}
\label{tab:chapter3-performance-metrics}
\begin{center}
\begin{tabular}{p{0.14\textwidth} p{0.28\textwidth} p{0.44\textwidth}}
\toprule
Metric & Description & Purpose \\
\midrule
\textbf{MAPE} & Mean Absolute Percentage Error & Primary metric; enables cross-study comparison \\
\textbf{R\textsuperscript{2}} & Coefficient of Determination & Proportion of variance explained \\
\textbf{MAE} & Mean Absolute Error (in PHP) & Business-interpretable error \\
\textbf{RMSE} & Root Mean Square Error & Penalizes large errors \\
\bottomrule
\end{tabular}
\end{center}
\end{table}

\subsubsection{\texorpdfstring{\textbf{Benchmark Studies Used for Later Comparison}}{Benchmark Studies Used for Later Comparison}}\label{benchmark-targets}

\begin{table}[htbp]
\caption{Benchmark Studies Used to Contextualize Later Evaluation Results}
\label{tab:chapter3-benchmark-studies}
\begin{center}
\begin{tabular}{p{0.22\textwidth} p{0.20\textwidth} p{0.10\textwidth} p{0.32\textwidth}}
\toprule
Study & Context & MAPE & Use in This Thesis \\
\midrule
Ramolete et al. (2023) & Philippines, larger dataset & 10--21\% & Provides a Philippine reference point for interpreting why cleaner, larger datasets can achieve lower percentage error. \\
Nyanda et al. (2024) & Tanzania, $n \approx 954$ & 48\% & Provides a small-sample emerging-market benchmark for the later evaluation chapter. \\
\bottomrule
\end{tabular}
\end{center}
\end{table}

\subsubsection{\texorpdfstring{\textbf{SHAP Explainability}}{SHAP Explainability}}\label{shap-explainability}

SHAP was used to interpret model outputs at both the dataset and property level, which was important to keep the results reviewable in valuation practice and consistent with IVS 2025 transparency requirements.

\begin{itemize}
\item \textbf{Global SHAP (Summary Plots)}: Identified which value drivers affected Metro Cebu property prices most across the entire dataset. This answered RQ1 (``What value drivers significantly influence property prices in Metro Cebu?'').
\item \textbf{Local SHAP (Force Plots)}: Explained individual predictions. For example: \emph{``This Mactan condominium is predicted higher because of stronger accessibility and island-location premiums, but the smaller floor area pulls the estimate down.''}
\end{itemize}

This helped keep the model interpretable enough to function as a support tool rather than a purely opaque prediction system. In practical terms, the model was evaluated not only by how accurately it predicted price, but also by whether its outputs could still be examined and discussed in terms that practitioners could understand.

\begin{center}\rule{0.5\linewidth}{0.5pt}\end{center}

\subsection{\texorpdfstring{\textbf{Deliverables}}{Deliverables}}\label{deliverables}

\subsubsection{\texorpdfstring{\textbf{QGIS Interactive Map}}{QGIS Interactive Map}}\label{qgis-interactive-map}

The main applied output of the study was a spatial review environment that organized model results across Metro Cebu. Rather than presenting results only as tables or summary statistics, the QGIS project was designed to show how predicted values, valuation-gap diagnostics, and locational patterns were distributed across the study area.

\begin{enumerate}
\item \textbf{Predicted Price Layers}: Point or surface outputs summarizing predicted open-market residential values across Metro Cebu.
\item \textbf{Valuation-Gap Diagnostics}: Comparison layers showing the divergence between model-based values and BIR zonal benchmarks.
\item \textbf{Supporting Locational Layers}: CBD-node, road-network, and amenity-accessibility context layers used to interpret the price surface geographically.
\end{enumerate}

The purpose of this output was to make the model easier to inspect geographically, especially for brokers, investors, and local government users who needed to compare values across locations rather than only at the level of individual records.

\subsubsection{\texorpdfstring{\textbf{Streamlit Web Application}}{Streamlit Web Application}}\label{streamlit-web-application}

A Streamlit web application complemented the QGIS environment by providing an interactive interface for inspecting the deployed open-market model. In practical use, the application centered on a predicted residential price surface for Metro Cebu while also supporting property-level prediction checks and SHAP-based explanation pages. While the QGIS output emphasized geographic comparison, the Streamlit application provided a more direct interface for exploring individual predictions and their underlying value drivers.

\begin{center}\rule{0.5\linewidth}{0.5pt}\end{center}
